
\documentclass{article}
\usepackage{colortbl}
\usepackage{makecell}
\usepackage{multirow}
\usepackage{supertabular}

\begin{document}

\newcounter{utterance}

\centering \large Interaction Transcript for game `chronicle', experiment `discoveries\_easy', episode 7 with llama{-}3.1{-}8b{-}instruct.
\vspace{24pt}

{ \footnotesize  \setcounter{utterance}{1}
\setlength{\tabcolsep}{0pt}
\begin{supertabular}{c@{$\;$}|p{.15\linewidth}@{}p{.15\linewidth}p{.15\linewidth}p{.15\linewidth}p{.15\linewidth}p{.15\linewidth}}
    \# & $\;$A & \multicolumn{4}{c}{Game Master} & $\;\:$B\\
    \hline

    \theutterance \stepcounter{utterance}  
    & & \multicolumn{4}{p{0.6\linewidth}}{
        \cellcolor[rgb]{0.9,0.9,0.9}{
            \makecell[{{p{\linewidth}}}]{
                \texttt{\tiny{[P1$\langle$GM]}}
                \texttt{You are playing a historical deduction game as the NARRATOR.} \\
\\ 
\texttt{You have been given the name of a historical event, discovery, or milestone. Your goal is to help the DETECTIVE figure out which event you are describing — but you must describe it indirectly, without ever naming it.} \\
\\ 
\texttt{Rules:} \\
\texttt{1. Each round, write exactly ONE sentence describing the event. Add NEW information each round — do not repeat what you wrote before.} \\
\texttt{2. You must NEVER use any of the following forbidden keywords (these would make the event too easy to identify): expedition, pancho, villa} \\
\texttt{3. You must NEVER state the event name directly or any close variant of it.} \\
\texttt{4. You may freely describe causes, consequences, the broader context, the sequence of events, and the significance — but only using general, non{-}identifying language.} \\
\\ 
\texttt{Format your response EXACTLY as follows (required):} \\
\texttt{PARAGRAPH: <exactly one sentence describing the event without forbidden keywords>} \\
\\ 
\texttt{The event you must describe (without naming it) is: Pancho Villa Expedition} \\
            }
        }
    }
    & \\ \\

    \theutterance \stepcounter{utterance}  
    & \multicolumn{4}{p{0.6\linewidth}}{
        \cellcolor[rgb]{0.9,0.9,0.9}{
            \makecell[{{p{\linewidth}}}]{
                \texttt{\tiny{[P1$\rangle$GM]}}
                \texttt{PARAGRAPH: In the early 20th century, a series of violent clashes between Mexican bandits and American soldiers along the US{-}Mexico border led to a significant escalation in tensions between the two nations.} \\
            }
        }
    }
    & & \\ \\

    \theutterance \stepcounter{utterance}  
    & & \multicolumn{4}{p{0.6\linewidth}}{
        \cellcolor[rgb]{0.9,0.9,0.9}{
            \makecell[{{p{\linewidth}}}]{
                \texttt{\tiny{[GM$\rangle$P2]}}
                \texttt{You are playing a historical deduction game as the DETECTIVE.} \\
\\ 
\texttt{A Narrator will describe a historical event, discovery, or milestone to you over several rounds. Each round you receive a new sentence of description that builds on what came before. Your job is to figure out which specific event is being described.} \\
\\ 
\texttt{After EACH new sentence, respond briefly. Keep your entire response under 100 words. Use this exact format (both sections are required):} \\
\\ 
\texttt{ANALYSIS: <2–3 short sentences: what does the growing description suggest? What type of event is this? What are the key features you are noticing?>} \\
\\ 
\texttt{GUESS: <your best guess for the name of the event being described — be specific, e.g. "The French Revolution" not just "a revolution">} \\
\\ 
\texttt{Rules:} \\
\texttt{{-} You may refine your GUESS each round as more information arrives.} \\
\texttt{{-} Your GUESS should name the specific event, not just a broad category.} \\
\texttt{{-} You have up to 10 rounds total.} \\
\texttt{{-} Be concise — quality of reasoning matters more than length.} \\
\\ 
\texttt{The first sentence from the Narrator is coming now.} \\
\\ 
\\ 
\texttt{{[}Sentence 1{]}} \\
\texttt{In the early 20th century, a series of violent clashes between Mexican bandits and American soldiers along the US{-}Mexico border led to a significant escalation in tensions between the two nations.} \\
            }
        }
    }
    & \\ \\

    \theutterance \stepcounter{utterance}  
    & & & \multicolumn{4}{p{0.6\linewidth}}{
        \cellcolor[rgb]{0.9,0.9,0.9}{
            \makecell[{{p{\linewidth}}}]{
                \texttt{\tiny{[GM$\langle$P2]}}
                \texttt{ANALYSIS: The description suggests a violent conflict between two groups, likely involving a border dispute or territorial issue. The involvement of bandits and soldiers implies a lawless or chaotic environment. The escalation of tensions between two nations indicates a significant international incident.} \\
\\ 
\texttt{GUESS: The Mexican Revolution} \\
            }
        }
    }
     \\ \\

    \theutterance \stepcounter{utterance}  
    & & \multicolumn{4}{p{0.6\linewidth}}{
        \cellcolor[rgb]{0.9,0.9,0.9}{
            \makecell[{{p{\linewidth}}}]{
                \texttt{\tiny{[P1$\langle$GM]}}
                \texttt{Round 1 complete. The Detective's current guess: "The Mexican Revolution". Please write a new paragraph with fresh information.} \\
            }
        }
    }
    & \\ \\

    \theutterance \stepcounter{utterance}  
    & \multicolumn{4}{p{0.6\linewidth}}{
        \cellcolor[rgb]{0.9,0.9,0.9}{
            \makecell[{{p{\linewidth}}}]{
                \texttt{\tiny{[P1$\rangle$GM]}}
                \texttt{PARAGRAPH: The event in question was sparked by a brazen attack on a small American town, which resulted in the deaths of several civilians and a significant amount of property damage, drawing widespread outrage and calls for retribution from the US government.} \\
            }
        }
    }
    & & \\ \\

    \theutterance \stepcounter{utterance}  
    & & \multicolumn{4}{p{0.6\linewidth}}{
        \cellcolor[rgb]{0.9,0.9,0.9}{
            \makecell[{{p{\linewidth}}}]{
                \texttt{\tiny{[GM$\rangle$P2]}}
                \texttt{{[}Sentence 2{]}} \\
\texttt{The event in question was sparked by a brazen attack on a small American town, which resulted in the deaths of several civilians and a significant amount of property damage, drawing widespread outrage and calls for retribution from the US government.} \\
            }
        }
    }
    & \\ \\

    \theutterance \stepcounter{utterance}  
    & & & \multicolumn{4}{p{0.6\linewidth}}{
        \cellcolor[rgb]{0.9,0.9,0.9}{
            \makecell[{{p{\linewidth}}}]{
                \texttt{\tiny{[GM$\langle$P2]}}
                \texttt{ANALYSIS: The description now suggests a specific, provocative incident that triggered a larger conflict. The attack on a civilian town and resulting property damage and civilian deaths indicate a brutal and indiscriminate act. The US government's response implies a strong reaction to the incident.} \\
\\ 
\texttt{GUESS: The Canutillo Raid} \\
            }
        }
    }
     \\ \\

    \theutterance \stepcounter{utterance}  
    & & \multicolumn{4}{p{0.6\linewidth}}{
        \cellcolor[rgb]{0.9,0.9,0.9}{
            \makecell[{{p{\linewidth}}}]{
                \texttt{\tiny{[P1$\langle$GM]}}
                \texttt{Round 2 complete. The Detective's current guess: "The Canutillo Raid". Please write a new paragraph with fresh information.} \\
            }
        }
    }
    & \\ \\

    \theutterance \stepcounter{utterance}  
    & \multicolumn{4}{p{0.6\linewidth}}{
        \cellcolor[rgb]{0.9,0.9,0.9}{
            \makecell[{{p{\linewidth}}}]{
                \texttt{\tiny{[P1$\rangle$GM]}}
                \texttt{PARAGRAPH: A high{-}ranking military officer, driven by a desire for revenge and a sense of national pride, led a large{-}scale military operation in response to the attack, which ultimately resulted in a series of bloody battles and a significant loss of life on both sides.} \\
            }
        }
    }
    & & \\ \\

    \theutterance \stepcounter{utterance}  
    & & \multicolumn{4}{p{0.6\linewidth}}{
        \cellcolor[rgb]{0.9,0.9,0.9}{
            \makecell[{{p{\linewidth}}}]{
                \texttt{\tiny{[GM$\rangle$P2]}}
                \texttt{{[}Sentence 3{]}} \\
\texttt{A high{-}ranking military officer, driven by a desire for revenge and a sense of national pride, led a large{-}scale military operation in response to the attack, which ultimately resulted in a series of bloody battles and a significant loss of life on both sides.} \\
            }
        }
    }
    & \\ \\

    \theutterance \stepcounter{utterance}  
    & & & \multicolumn{4}{p{0.6\linewidth}}{
        \cellcolor[rgb]{0.9,0.9,0.9}{
            \makecell[{{p{\linewidth}}}]{
                \texttt{\tiny{[GM$\langle$P2]}}
                \texttt{ANALYSIS: The description now suggests a large{-}scale military operation led by a high{-}ranking officer, driven by personal and national motivations. The series of bloody battles and significant loss of life on both sides indicate a prolonged and intense conflict. The officer's desire for revenge and national pride implies a personal stake in the conflict.} \\
\\ 
\texttt{GUESS: The Punitive Expedition} \\
            }
        }
    }
     \\ \\

    \theutterance \stepcounter{utterance}  
    & & & \multicolumn{2}{p{0.3\linewidth}}{
        \cellcolor[rgb]{0.9,0.9,0.9}{
            \makecell[{{p{\linewidth}}}]{
                \texttt{\tiny{[GM$|$GM]}}
                \texttt{round 3} \\
            }
        }
    }
    & & \\ \\

\end{supertabular}
}

\end{document}
